\uuid{XFWM}
\exo7id{3933}
\titre{exo7 3933}
\auteur{quercia}
\organisation{exo7}
\datecreate{2010-03-11}
\isIndication{false}
\isCorrection{false}
\chapitre{Dérivabilité des fonctions réelles}
\sousChapitre{Autre}
\cpge{mpsi-mp2i.fonctions-valeurs-reelles-complexes.a, pcsi.techniques-calcul-fonctions.a, ptsi.techniques-calcul-fonctions.a, mpsi-mp2i.fonctions-valeurs-reelles-complexes.b, pcsi.techniques-calcul-fonctions.b, ptsi.techniques-calcul-fonctions.b, bcpst1.fonctions-reelles-usuelles.a, bcpst1.derivees-primitives-integrales.a}
\module{Analyse}
\niveau{L1}
\difficulte{}

\contenu{
\texte{
Soit $f:\R \to \R$ dérivable.
}
\begin{enumerate}
    \item \question{Que peut-on dire de $f'$ si on sait que $f$ est paire ? impaire ? périodique ?}
    \item \question{Que peut-on dire de $f$ si on sait que $f'$ est paire ? impaire ? périodique ?}
    \item \question{Montrer que si $f'$ est $T$-périodique et $f(T) \ne f(0)$, alors $f$ n'a pas de
    période (on étudiera $f(nT)$ pour $n\in\N$).}
\end{enumerate}
}
